\documentclass[12pt]{article}
\usepackage{amsmath, amssymb}
\usepackage[margin=1in]{geometry}
\setlength{\parskip}{6pt}
\title{QUBO Formulation: Benchmark Job Selection Problem}
\date{}
\begin{document}
\maketitle

\subsection*{Benchmark Job Selection Problem}

\paragraph{Problem Statement.}
Given a set of $N$ candidate benchmark jobs, each characterized by a coverage reward $v_i$, a runtime $t_i$, and pairwise redundancy penalties $p_{ij}$, select a subset of jobs to maximize the total coverage reward while penalizing redundant coverage, subject to a total compute-time budget $T$.

\paragraph{Decision Variables.}
For each candidate job $i \in \{0, 1, \dots, N-1\}$, introduce a binary decision variable $x_i \in \{0,1\}$, where $x_i = 1$ if job $i$ is selected, and $x_i = 0$ otherwise.

\paragraph{QUBO Derivation.}
Step 1 --- The original objective is a maximization problem. Standard QUBO solvers minimize energy, so we negate the objective:
\begin{align}
\min_{x \in \{0,1\}^N} \left( -\sum_{i=0}^{N-1} v_i x_i + \sum_{0 \le i < j \le N-1} p_{ij} x_i x_j \right).
\end{align}

Step 2 --- The compute-time budget constraint is expressed as:
\begin{align}
\sum_{i=0}^{N-1} t_i x_i \le T.
\end{align}

Step 3 --- We convert the hard constraint into a quadratic penalty term. Following the standard QUBO convention, we add $\lambda \left( \sum_{i=0}^{N-1} t_i x_i - T \right)^2$ to the objective, where $\lambda > 0$ is a penalty weight. Expanding the square yields:
\begin{align}
\lambda \left( \sum_{i=0}^{N-1} t_i^2 x_i^2 - 2T \sum_{i=0}^{N-1} t_i x_i + T^2 + \sum_{0 \le i < j \le N-1} 2 t_i t_j x_i x_j \right).
\end{align}
Using the idempotent property of binary variables, $x_i^2 = x_i$, the expression simplifies to:
\begin{align}
\lambda \sum_{i=0}^{N-1} (t_i^2 - 2T t_i) x_i + \lambda \sum_{0 \le i < j \le N-1} 2 t_i t_j x_i x_j + \lambda T^2.
\end{align}

Step 4 --- Combining the negated objective and the expanded penalty term gives the full QUBO Hamiltonian $H(x)$:
\begin{align}
H(x) = \sum_{i=0}^{N-1} \left[ -v_i + \lambda(t_i^2 - 2T t_i) \right] x_i + \sum_{0 \le i < j \le N-1} \left[ p_{ij} + 2\lambda t_i t_j \right] x_i x_j + \lambda T^2.
\end{align}

Step 5 --- Reading off the Q matrix entries and the constant offset $c$ from $H(x)$ in general closed form:
\begin{align}
Q_{i,i} &= -v_i + \lambda(t_i^2 - 2T t_i), \quad i \in \{0, \dots, N-1\}, \\
Q_{i,j} &= p_{ij} + 2\lambda t_i t_j, \quad 0 \le i < j \le N-1, \\
c &= \lambda T^2.
\end{align}

\paragraph{Penalty Weight Justification.}
To guarantee that any budget violation is strictly penalized more than the maximum possible gain in the objective function, we require:
\begin{align}
\lambda > \frac{\sum_{i=0}^{N-1} v_i}{(\min_{i} t_i)^2}.
\end{align}
This condition ensures that the quadratic penalty incurred by exceeding the budget dominates the total reward achievable by any feasible subset, thereby enforcing feasibility at the global minimum.

\paragraph{Correctness Argument.}
Minimizing $H(x)$ over $\{0,1\}^N$ recovers the optimal solution because (i) all feasible assignments satisfy $\sum_i t_i x_i \le T$, yielding a penalty term that evaluates to a constant offset $c$ (or zero relative slack), and (ii) any infeasible assignment violates the budget by at least $\min_i t_i$, incurring a penalty of at least $\lambda (\min_i t_i)^2$. By the choice of $\lambda$, this penalty strictly exceeds the maximum possible reduction in the objective function from any feasible configuration. Consequently, the global minimum of $H(x)$ must correspond to a feasible assignment that optimally balances coverage rewards and redundancy penalties.

\paragraph{Illustrative Example.}
Consider a concrete instance with $N=3$ jobs. Let the parameters be:
\begin{align}
v_0 = 10, \; v_1 = 8, \; v_2 = 12; \quad t_0 = 3, \; t_1 = 4, \; t_2 = 5; \quad T = 7; \quad p_{01} = 2, \; p_{02} = 3, \; p_{12} = 1.
\end{align}
Choose $\lambda = 10$. Substituting these values into the general formulas yields the concrete Hamiltonian:
\begin{align}
H(x) &= -340 x_0 - 408 x_1 - 462 x_2 + 242 x_0 x_1 + 303 x_0 x_2 + 401 x_1 x_2 + 490.
\end{align}
Evaluating $H(x)$ over all feasible subsets ($\sum_i t_i x_i \le 7$):
\begin{align}
H(\emptyset) &= 490, \quad H(\{0\}) = 150, \quad H(\{1\}) = 82, \quad H(\{2\}) = 28, \quad H(\{0,1\}) = -16.
\end{align}
The global minimum is $H(\{0,1\}) = -16$, corresponding to the optimal selection of jobs $0$ and $1$, which yields a net coverage value of $v_0 + v_1 - p_{01} = 16$ while respecting the runtime budget $t_0 + t_1 = 7 \le T$.
\end{document}